\documentclass[../../main.tex]{subfiles}
\begin{document}
\begin{oframed}
  \subfile{problem}
\end{oframed}

{\bf [解答]}

\begin{figure}[htb]
  \centering
  \begin{tikzpicture}[scale=1.4, >=stealth, every node/.style={font=\small}]
    % 座標軸
    \draw[->, thick, gray!80!black] (-3.2, 0) -- (3.5, 0) node[right] {$x$};
    \draw[->, thick, gray!80!black] (0, -2.5) -- (0, 2.8) node[above] {$y$};
    \node[below left=2pt] at (0, 0) {$O$};

    % 双曲線 xy = 1 (第1象限および第3象限)
    \draw[line width=1.2pt, blue!85!black, domain=0.38:3.2, samples=70]
    plot (\x, {1.0 / \x}) node[right] {$xy=1$};
    \draw[line width=1.2pt, blue!85!black, domain=-3.2:-0.38, samples=70]
    plot (\x, {1.0 / \x});

    % 代表点 P (第1象限) と Q (第3象限)
    % a = 3^{3/4} \approx 2.28, c = 3^{-1/4} \approx 0.76
    \def\aX{2.28}
    \def\aY{0.439}
    \def\cX{-0.76}
    \def\cY{-1.316}

    \coordinate (P) at (\aX, \aY);
    \coordinate (Q) at (\cX, \cY);

    % 線分 QP
    \draw[line width=1.3pt, red!85!black] (Q) -- (P);

    % 水平補助線と角度 30^\circ
    \draw[dashed, gray!70!black] (Q) -- ++(1.8, 0);
    \draw[->, thick, orange!85!black] ($(Q) + (0.8, 0)$) arc (0:30:0.8);
    \node[orange!85!black, font=\footnotesize] at ($(Q) + (15:1.1)$) {$30^\circ$};

    % 点ノード
    \node[circle, fill=blue!85!black, inner sep=1.3pt] at (P) {};
    \node[above right=1pt] at (P) {$P\left(a, \frac{1}{a}\right)$};

    \node[circle, fill=blue!85!black, inner sep=1.3pt] at (Q) {};
    \node[below left=1pt] at (Q) {$Q\left(-c, -\frac{1}{c}\right)$};

    % 投影点線
    \draw[densely dotted, thick, gray!70!black] (P) -- (\aX, 0) node[below] {$a$};
    \draw[densely dotted, thick, gray!70!black] (P) -- (0, \aY) node[left] {$\frac{1}{a}$};
    \draw[densely dotted, thick, gray!70!black] (Q) -- (\cX, 0) node[above] {$-c$};
    \draw[densely dotted, thick, gray!70!black] (Q) -- (0, \cY) node[right] {$-\frac{1}{c}$};
  \end{tikzpicture}
  \caption{双曲線 $xy=1$ 上の定点 $P$ と動点 $Q$ および線分 $QP$}
  \label{fig:hyperbola_qp}
\end{figure}

題意から
\begin{align}
  a>0, \quad
  b=\dfrac{1}{a}
\end{align}
である．また$c>0$として$\left(-c,\dfrac{-1}{c}\right)$とおける．

\paragraph{(1)}

$|QP|^2=(c+a)^2+\left(\dfrac{1}{c}+\dfrac{1}{a}\right)^2$であるからこれを$f(c)$として
\begin{align}
  \dfrac{1}{2}f'(c)&=(c+a)+\left(\dfrac{1}{c}+\dfrac{1}{a}\right)\left(\dfrac{-1}{c^2}\right)  \\
  &=(a+c)\left(1-\dfrac{1}{ac^3}\right)  \\
  &=\dfrac{a+c}{ac^3}{(pc-1)\left(1+pc+(pc)^2\right)}
\end{align}
となる．ただし$p=\sqrt[3]{a}$である．故に下表を得る．

\begin{align}
  \begin{array}{|c|c|c|c|} \hline
    c  &   &  1/p   &      \\ \hline
    f'  & - &   0     & +     \\ \hline
    f  &\searrow&         &\nearrow  \\ \hline
  \end{array}
\end{align}
従って
\begin{align}
  \min f(c)&=f\left(\dfrac{1}{p}\right) \\
  &=\left(\dfrac{1}{p}+p^3\right)^2+\left(p+\dfrac{1}{p^3}\right)^2  \\
  &=\left(p^2+\dfrac{1}{p^2}\right)^3
\end{align}
だから，もとめる最小値は($|QP|\ge0$より)
\begin{align}
  \min|QP|=(a^{2/3}+a^{-2/3})^{3/2} \tag{答}
\end{align}
となる．

\paragraph{(2)}
このとき，前問の結果から，
\begin{align}
  \overrightarrow{QP}&=
  \begin{pmatrix}p^3+1/p\\1/p^3+p
  \end{pmatrix}  \\
  &\parallel
  \begin{pmatrix}p^2\\1
  \end{pmatrix}
\end{align}
となる．従ってこれが$x$軸の正の方向と$30^\circ$の角を成す，つまり$(\sqrt{3},1)$と平行なとき
\begin{align}
  p^2&=\sqrt{3}　\\
  p&=\sqrt[4]{3} \\
  a&=3^{3/4}\tag{答}
\end{align}
となる．
\newpage
\end{document}
